\section{Discusión}
\label{sec:discusion}

\subsection{Contraste de la hipótesis con las pruebas en Melanir}
La hipótesis planteó una mediana de RTT menor con sockets TCP para
solicitudes cortas y secuenciales en localhost, usando conexiones ya
establecidas. En la sesión de Melanir, sockets obtuvo una mediana de
0,0901 ms frente a 1,28585 ms de gRPC con solicitudes serializadas de
100 bytes. La mediana de sockets fue un 92,99\,\% menor, calculado
respecto de gRPC. Esta observación respalda la dirección de la hipótesis
en las condiciones ejecutadas. La hipótesis se refería aproximadamente
al contenido del mensaje, mientras la implementación controla el tamaño
serializado completo; esa diferencia acota su comprobación.

Los porcentajes de esta discusión se calculan con los valores completos
de \texttt{resumen.csv}, antes del redondeo de las tablas. Para la
reducción relativa se utiliza
\[
 100\left(1-\frac{\text{latencia de sockets}}{\text{latencia de gRPC}}\right).
\]
Las medias secuenciales fueron 0,1136 ms y 1,3079 ms con 100 bytes,
y 0,1644 ms y 1,4494 ms con 10\,000 bytes, para sockets y gRPC,
respectivamente. Las reducciones de la media fueron 91,31\,\% y
88,66\,\%. En concurrencia también se observó menor media en sockets:
86,21\,\% menos con 100 bytes y 80,28\,\% menos con 10\,000 bytes.
Estos porcentajes describen esta sesión; no constituyen una ventaja
universal de sockets sobre gRPC.

El coste de las bibliotecas y del tratamiento de la llamada es una
explicación compatible con el resultado, pero no se midió por separado.
El JSON se construye antes de iniciar el cronómetro de sockets; la
serialización interna de la llamada gRPC queda dentro de su RTT.
Tampoco se alternó el orden de ejecución de los mecanismos.
Por tanto, no podemos asignar el porcentaje observado exclusivamente
a HTTP/2, a Protocol Buffers o a TCP. Para separar esos efectos haría
falta instrumentar cada etapa y repetir con orden alternado.

\subsection{Rendimiento y concurrencia}
Con tres clientes y 100 bytes, las tasas agregadas fueron
11\,836,00 mensajes/s en sockets y 1\,387,18 mensajes/s en gRPC:
una relación de 8,53 a 1. Con 10\,000 bytes, esa relación bajó a
4,89 a 1, con 6\,428,64 y 1\,315,57 mensajes/s. El tamaño de la
solicitud cambió la separación entre las implementaciones, de modo que
una única carga no habría bastado para describirlas.

La tasa concurrente superó la secuencial en ambos mecanismos, aunque
también aumentó la latencia media por solicitud. Por ejemplo, sockets
con 100 bytes pasó de 0,1136 a 0,1866 ms. Enviar desde tres clientes
permite solapar trabajo y completar más mensajes en conjunto, sin
asegurar que cada uno termine antes.
El tiempo total incluye preparación y cierre de clientes y conexiones,
y el grupo de hilos en concurrencia; las tasas comparan ese ciclo
completo, no solo solicitudes sobre conexiones preparadas.
Las barras de error muestran variación entre cinco medias de corrida.
No se efectuaron contrastes de significación ni se midió el punto de
saturación del servidor.

\subsection{Desarrollo, tipado y manejo de errores}
La implementación TCP del repositorio permite ver cada paso: abrir
el socket, construir el JSON y delimitarlo mediante un salto de línea.
Para este intercambio pequeño, el flujo resulta directo de seguir.
El equipo debe conservar por su cuenta el acuerdo sobre campos,
delimitación y validación. En \texttt{src/sockets/nodo.py}, la recepción
interpreta el JSON y accede a los campos sin una validación previa
del esquema; una entrada mal formada puede provocar una excepción.

gRPC requiere mantener el contrato \texttt{.proto} y generar módulos
antes de ejecutar los nodos. A cambio, los tipos de solicitud y respuesta
están declarados y el código utiliza la operación
\texttt{ExchangeMessage}. El esquema diferencia la hora física textual
del reloj lógico entero. Ese tipado no valida todo el significado:
el servidor todavía comprueba el emisor y el separador del identificador,
y devuelve \texttt{INVALID\_ARGUMENT} cuando no son válidos.

El cliente gRPC configura una espera de disponibilidad de 10 segundos
y un plazo de 5 segundos para la llamada de intercambio. Comprueba
también el contenido del acuse. El nodo TCP publicado no ofrece un
acuse de aplicación en su intercambio de eventos ni configura esos
plazos explícitos. No debe confundirse ese intercambio con los
servidores de eco utilizados para medir RTT. Estas diferencias afectan
a la funcionalidad y a la forma de detectar fallos, aunque la sesión
de latencia no ensayó desconexiones ni solicitudes inválidas.
Por eso no se puede concluir que una implementación sea más robusta
a partir de sus tiempos normales.

La verificación del reloj respalda las relaciones registradas, pero
no mide el tiempo físico entre eventos. En gRPC, el acuse de la llamada
no se registra como evento adicional ni actualiza el reloj del cliente
al llegar. Revisar ese recorrido sigue siendo necesario para modelar
el ciclo completo de solicitud y respuesta.

\subsection{Aplicación al PFC y límites de la decisión}
El PFC del equipo de Angel es el SGIRI, un sistema de gestión de
incidencias de red para un proveedor de internet, que se está separando
en servicios. En él, gRPC encaja en la consulta de
\texttt{servicio-incidencias} a \texttt{servicio-tecnicos} para
resolver el técnico asignado a una incidencia: es una operación entre
servicios con solicitud y respuesta definidas, donde el contrato
\texttt{.proto} y el plazo de la llamada permiten responder con un
error controlado si el otro servicio no contesta. Sockets TCP encaja
en un canal interno pequeño, como el envío periódico del estado de un
equipo de red hacia el servicio de incidencias, donde el formato es
simple y lo mantiene el propio equipo.

En las mediciones del 22 de noviembre en el equipo Mourad, la mediana
de gRPC estuvo entre 0,34 y 0,58~ms y la de sockets entre 0,05 y
0,15~ms. La diferencia es menor a un milisegundo en loopback, por lo
que en el SGIRI la elección depende más del contrato y del manejo de
errores que de la latencia. Esta decisión se documentará en el ADR del
mecanismo de comunicación del PFC.

Las mediciones en loopback dejan fuera la latencia de una red entre
máquinas y comparten los recursos del mismo equipo. Como detalla
el apartado~\ref{sec:validez}, también quedan diferencias de
serialización y carga de fondo. La elección final debe contrastar
latencia, contrato, mantenimiento y comportamiento ante fallos en el
escenario real. Las fuentes de RPC y serialización sitúan esa comparación
en un marco más amplio \cite{blinowski2026rpc,jackson2024streaming};
las cifras anteriores proceden únicamente de nuestros CSV.
